% backmatter/supplement/s-ch4-summary.tex -- "Summary: Corrections and Additions to Altitude Equations for 2-Stage
% and Multistage Rockets" (PDF 703-710; typewritten with handwritten displays; no author or date shown).
% Every display is unnumbered, as on the page. Notation follows STYLE.md section 15 (Chapter 4); the upright
% subscripts term, fterm and BC are word subscripts; v_{(n-1)} of the rewritten (50)-(51) keeps its printed
% parentheses. Equation citations without a page number are \eqref's to the chapter; "page 531" etc. stay literal.
% A few short displays that follow a short text line are gather* (full display skip) rather than equation*.
\chapter{Summary: Corrections and Additions to Altitude Equations for 2-Stage and Multistage Rockets}
\label{supp:ch4-summary}
\markright{\MakeUppercase{Summary: Corrections and Additions to Altitude Equations}}% the full title overruns the running head

\noindent\edcap{The equation rewrites of this Summary, with the definitions of $\hat{t}$, $\hat{y}$ and
$I_{t2}$ and the change to the discussion of the upper limit, are applied in Section~\ref{ch4:sec:2.2} of
Chapter~\ref{ch4}; its plans for new material are not carried out there (see the editor's note in
Section~\ref{ch4:sec:2.2.1}).}

\phantomsection
\section*{Second-Stage Burning-Phase Solutions}\label{supp:ch4-summary:burning}

In hindsight, better physical insight into the correct limits of integration for solving the multistage
altitude performance equations could have been afforded by formulating the upper-stage burning-phase
equations as an initial-value problem. At $t = t_1$, when the second-stage engine ignites, the rocket is
traveling at an ``initial velocity'' of $v_1$. If we define the time elapsed since the ignition of the
second-stage engine as $\hat{t} = (t - t_1)$ and the additional altitude gained since the ignition of the
second-stage engine as $\hat{y} = (y - y_1)$, we can summarize the changes in the rocket's velocity and
altitude that occur during the second-stage burning time as follows:
\begin{itemize}
  \item $v$ changes from $v_1$ to $v_2$ as $\hat{t}$ increases from 0 to $t_2$.
  \item $\hat{y}$ increases from 0 to $y_2$ as $\hat{t}$ increases from 0 to $t_2$.
\end{itemize}

This formulation automatically supplies the correct limits of integration for both the Caporaso-Bengen and
the Fehskens-Malewicki solutions, and it simplifies the second-stage equations by expressing them in a form
that is identical to that of the single-stage equations except for the existence of a nonzero initial
velocity in the second-stage equations.

I plan to move the notation used for two-staged rockets from Section~\ref{ch4:sec:2.2.1} to
Section~\ref{ch4:sec:2.2}, since it is used for both the Fehskens-Malewicki and the Caporaso-Bengen
solution, and to add the following variables:
\begin{itemize}[nosep,leftmargin=4em,labelwidth=3em,labelsep=0.35em,align=right]% hanging indent, as printed
  \item[$v$ =] velocity at any time during the burning of the second-stage engine
  \item[$\hat{y}$ =] altitude gained since the ignition of the second-stage engine (not the total altitude
  above the launch point)
  \item[$\hat{t}$ =] time elapsed since the ignition of the second-stage engine
\end{itemize}

I also plan to move the discussion of the physics of two-staged flight from~\ref{ch4:sec:2.2.1}
to~\ref{ch4:sec:2.2}, since it applies to all methods of solution.

\phantomsection
\section*{Extended Caporaso-Bengen Solution}\label{supp:ch4-summary:cb}

Rewrite equation~\eqref{ch4:eq:40} on page 531 of the original text as
\begin{equation*}
\int_{0}^{t_2} m_2\frac{dv}{d\hat{t}}\,d\hat{t} = \int_{0}^{t_2} F_2(\hat{t})\,d\hat{t}
  - \int_{0}^{t_2} m_2g\,d\hat{t} - \int_{0}^{t_2} k_2v^{2}\,d\hat{t}
\end{equation*}
This makes the discussion of the additive properties of integrals taken between adjacent limits, and the
later discussion of treating $k$ as if it had been $k_2$ throughout the flight, unnecessary.

Rewrite equation~\eqref{ch4:eq:44} on page 532 of the original text as
\begin{equation*}
\int_{0}^{t_2} k_2v\frac{d\hat{y}}{d\hat{t}}\,d\hat{t} = k_2\hat{y}v\Big]_{0,v_1}^{y_2,v_2}
  - k_2\int_{0}^{t_2}\hat{y}\frac{dv}{d\hat{t}}\,d\hat{t} \cong k_2y_2v_2
\end{equation*}
This is a correction. The original evaluation of the truncated drag integral was wrong. The second stage's
momentum is changed during the burning time of the second-stage engine only by the forces acting on it
during that time. At $\hat{t} = 0$ $v = v_1$ but $\hat{y} = 0$. Because the increment of altitude gained
since second-stage ignition is zero at the moment of ignition, zero should have appeared instead of $y_1$ in
the original presentation of~\eqref{ch4:eq:44}. Our physical intuition tells us that $y_1$ should not
affect the value of $v_2$ or $y_2$, since the equations as presented do not explicitly account for the
reduction of atmospheric density with increasing altitude, and the correctly-truncated drag integral
confirms that this is indeed the case. Correcting equation~\eqref{ch4:eq:44} causes all terms in $y_1$ to
disappear from the subsequent equations and thus simplifies their algebra.

Corrected equation~\eqref{ch4:eq:45} of the original text becomes
\begin{equation*}
m_2v_2 - m_2v_1 = I_{t2} - m_2gt_2 - k_2y_2v_2
\end{equation*}
where I have substituted $I_{t2}$, the total impulse of the second-stage engine, instead of $F_2t_2$ for
$\int_{0}^{t_2} F_2(\hat{t})\,d\hat{t}$. As in the single-stage case, the burnout velocity does not depend on
the exact functional form of the thrust, but the burnout altitude does. Corrected
equation~\eqref{ch4:eq:46} then becomes
\begin{equation*}
v_2 = \frac{I_{t2} - m_2gt_2 + m_2v_1}{m_2 + k_2y_2}
\end{equation*}
When $F_2(\hat{t})$ is approximated by the constant, average value $F_2$ as in the single-stage case and
corrected equation~\eqref{ch4:eq:45} is integrated, corrected equation~\eqref{ch4:eq:47} becomes
\begin{equation*}
m_2y_2 - m_2v_1t_2 = \frac{t_2^{2}}{2}(F_2 - m_2g) - k_2\frac{y_2^{2}}{2}
\end{equation*}
This equation was actually evaluated correctly in the original text, except for the terms in $y_1$. The
time-dependent variables were all assumed to be functions of $(t - t_1)$ and the integration was performed
between 0 and the variable upper limit $(t - t_1)$. Once the functional forms of the integrals had been thus
determined, they were evaluated between a lower limit of 0 and an upper limit of $[(t_1 + t_2) - t_1] =
t_2$. The approach was essentially the same as the one I propose to use for the rewrite, except that no new
name was assigned to the quantity $(t - t_1)$.

Rewrite corrected equation~\eqref{ch4:eq:48} on page 533 of the original text as
\begin{equation*}
y_2 = \frac{-m_2 + \sqrt{m_2^{2} + 2k_2t_2\bigl[\tfrac{t_2}{2}(F_2 - m_2g) + m_2v_1\bigr]}}{k_2}
\end{equation*}
This introduces another correction. The original equation has an algebraic sign error in the last term of
the square root. The quantity shown as $-m_2v_1$ should be $+m_2v_1$. Larry Curcio saw this error before I
did, as his marked xerox copy of page 533 shows. I only discovered it when I rederived the equation.

Rewrite corrected equation~\eqref{ch4:eq:49} as
\begin{equation*}
v_2 = \frac{I_{t2} - m_2gt_2 + m_2v_1}{\sqrt{m_2^{2} + 2k_2t_2\bigl[\tfrac{t_2}{2}(F_2 - m_2g) + m_2v_1\bigr]}}
\end{equation*}
Revise the discussion of limiting behavior to eliminate the reference to $y_1$. Rewrite corrected
equation~\eqref{ch4:eq:50} as
\begin{equation*}
v_n = \frac{I_{tn} - m_ngt_n + m_nv_{(n-1)}}{\sqrt{m_n^{2} + 2k_nt_n\bigl[\tfrac{t_n}{2}(F_n - m_ng)
  + m_nv_{(n-1)}\bigr]}}
\end{equation*}
and rewrite corrected equation~\eqref{ch4:eq:51} as
\begin{equation*}
y_n = \frac{-m_n + \sqrt{m_n^{2} + 2k_nt_n\bigl[\tfrac{t_n}{2}(F_n - m_ng) + m_nv_{(n-1)}\bigr]}}{k_n}
\end{equation*}

\phantomsection
\section*{Extended Fehskens-Malewicki Solution}\label{supp:ch4-summary:fm}

Rewrite equation~\eqref{ch4:eq:52} as
\begin{gather*}
\int_{v_1}^{v_2}\frac{m_2\,dv}{F_2 - m_2g - k_2v^{2}} = \int_{0}^{t_2} d\hat{t}
\end{gather*}
and rewrite equation~\eqref{ch4:eq:53} as
\begin{equation*}
\frac{m_2}{\sqrt{k_2(F_2 - m_2g)}}\tanh^{-1}\biggl[v\sqrt{\frac{k_2}{F_2 - m_2g}}\,\biggr]_{v_1}^{v_2}
  = t_2
\end{equation*}
Change discussion of upper limit $t$ to upper limit $\hat{t}$, where $\hat{t}$ is greater than or equal to 0
and less than or equal to $t_2$.

Move the discussion of the generalization of the equations to rockets of three or more stages to the end of
the new section on qualitative features, limiting behavior, and terminal velocity.

\phantomsection
\section*{Qualitative Features and Limiting Behavior of Multistage Solutions.\\ Considerations of Terminal
Velocity}\label{supp:ch4-summary:terminal}

Add a Section 2.2.3 which discusses qualitative features and limiting behavior of multistage solutions, and
treats the subject of approach to terminal velocity from below and from above. State that attainment of
terminal velocity
\begin{gather*}
v_{\mathrm{term}} = \sqrt{\frac{F_2 - m_2g}{k_2}}
\end{gather*}
requires the assumption of constant average thrust. Show that, when corrected equation~\eqref{ch4:eq:49} is
modified to incorporate that assumption and is written as
\begin{equation*}
v_2 = \frac{t_2(F_2 - m_2g) + m_2v_1}{\sqrt{m_2^{2} + 2k_2t_2\bigl[\tfrac{t_2}{2}(F_2 - m_2g) + m_2v_1\bigr]}}
\end{equation*}
then when $v_1 = v_{\mathrm{term}}$, the substitution of $k_2v_1^{2}$ for $(F_2 - m_2g)$
in~\eqref{ch4:eq:48} and~\eqref{ch4:eq:49} simplifies them to
\begin{equation*}
y_2 = v_1t_2
\end{equation*}
and
\begin{equation*}
v_2 = v_1
\end{equation*}
Note that the extended Caporaso-Bengen solution predicts terminal velocity exactly, because truncating the
drag integral causes no error when the velocity is constant $[(dv/dt) = 0]$.

Show that the extended Fehskens-Malewicki solution can also be written in terms of $v_{\mathrm{term}}$ as
\begin{equation*}
v_2 = (v_{\mathrm{term}})\tanh\biggl[\frac{k_2t_2}{m_2}v_{\mathrm{term}}
  + \tanh^{-1}\biggl(\frac{v_1}{v_{\mathrm{term}}}\biggr)\biggr]
\end{equation*}
and
\begin{equation*}
y_2 = \frac{m_2}{k_2}\ln\biggl[\cosh\biggl(\frac{k_2t_2}{m_2}v_{\mathrm{term}}\biggr)
  + \frac{v_1}{v_{\mathrm{term}}}\sinh\biggl(\frac{k_2t_2}{m_2}v_{\mathrm{term}}\biggr)\biggr]
\end{equation*}
As $v_1$ approaches $v_{\mathrm{term}}$ from below, the inverse hyperbolic tangent in the velocity equation
increases without bound and dominates the quantity in brackets:
\begin{equation*}
\frac{k_2t_2}{m_2}v_{\mathrm{term}} + \tanh^{-1}\biggl(\frac{v_1}{v_{\mathrm{term}}}\biggr)
  \quad\longrightarrow\quad \tanh^{-1}\biggl(\frac{v_1}{v_{\mathrm{term}}}\biggr)
\end{equation*}
so that
\begin{equation*}
v_2 \longrightarrow (v_{\mathrm{term}})\tanh\biggl[\tanh^{-1}\biggl(\frac{v_1}{v_{\mathrm{term}}}\biggr)\biggr]
\end{equation*}
and $v_2$ approaches $v_1$. Direct substitution of $v_1$ for $v_{\mathrm{term}}$ in the altitude increment
equation gives the result
\begin{equation*}
y_2 = \frac{m_2}{k_2}\ln\biggl[\cosh\biggl(\frac{k_2t_2}{m_2}v_1\biggr)
  + \sinh\biggl(\frac{k_2t_2}{m_2}v_1\biggr)\biggr]
\end{equation*}
from which, using the exponential definitions of $\cosh$ and $\sinh$, we obtain
\begin{equation*}
y_2 = \frac{m_2}{k_2}\ln\Bigl[e^{\bigl(\frac{k_2t_2}{m_2}v_1\bigr)}\Bigr]
\end{equation*}
from which
\begin{equation*}
y_2 = v_1t_2
\end{equation*}
The extended Fehskens-Malewicki solution therefore also predicts terminal velocity exactly, despite the
singularity in the inverse hyperbolic tangent at $(v_1/v_{\mathrm{term}}) = 1$.

Describe the case of the low-thrust sustainer and approach to terminal velocity from above. Explain that the
extended Caporaso-Bengen solution has no singularity in either its velocity equation or its altitude
increment equation at terminal velocity, and that the equations have the same functional forms when $v_1$ is
greater than $v_{\mathrm{term}}$ that they do when $v_1$ is less than $v_{\mathrm{term}}$. Discuss
degradation in accuracy of extended Caporaso-Bengen solution when $v_1$ is much greater than
$v_{\mathrm{term}}$ due to inaccuracy of truncated drag integral when $dv/dt$ is of large magnitude.
Numerical experimentation to date indicates $v_1$ must not be greater than $\sqrt{2}\cdot(v_{\mathrm{term}})$
for extended Caporaso-Bengen solution to produce acceptably accurate results.

Note that the extended Fehskens-Malewicki velocity equation (state equation numbers) derived on the
assumption that $v_1$ is less than $v_{\mathrm{term}}$ cannot be used when $v_1$ is greater than
$v_{\mathrm{term}}$ because the inverse hyperbolic tangent is undefined for an argument greater than 1.
Explain that the integral
\begin{equation*}
\int_{v_1}^{v_2}\frac{m_2\,dv}{F_2 - m_2g - k_2v^{2}}
\end{equation*}
of the extended Fehskens-Malewicki velocity equation has a different functional form when the quantity
\begin{equation*}
F_2 - m_2g - k_2v^{2}
\end{equation*}
is less than zero. In such cases the integral is
\begin{equation*}
\frac{m_2}{\sqrt{k_2(F_2 - m_2g)}}\coth^{-1}\biggl[v\sqrt{\frac{k_2}{F_2 - m_2g}}\,\biggr]_{v_1}^{v_2}
\end{equation*}
from which the velocity is obtained as
\begin{equation*}
v_2 = \sqrt{\frac{F_2 - m_2g}{k_2}}\,\coth\biggl[\tfrac{t_2}{m_2}\sqrt{k_2(F_2 - m_2g)}
  + \coth^{-1}\biggl(v_1\sqrt{\frac{k_2}{F_2 - m_2g}}\biggr)\biggr]
\end{equation*}
The altitude increment gained during the second-stage burn can also be determined from this equation, by
integrating it from $\hat{t} = 0$ to $\hat{t} = t_2$:
\begin{equation*}
y_2 = \frac{m_2}{k_2}\ln\left\{\frac{\sinh\Bigl[\tfrac{t_2}{m_2}\sqrt{k_2(F_2 - m_2g)}
  + \coth^{-1}\Bigl(v_1\sqrt{\frac{k_2}{F_2 - m_2g}}\Bigr)\Bigr]}{\sinh\Bigl[\coth^{-1}\Bigl(v_1
  \sqrt{\frac{k_2}{F_2 - m_2g}}\Bigr)\Bigr]}\right\}
\end{equation*}
and through the use of the well-known identity
\begin{equation*}
\sinh(A + B) = \sinh(A)\cosh(B) + \cosh(A)\sinh(B)
\end{equation*}
where $A$ and $B$ represent any two quantities, the altitude equation may be cast into the simpler form
\begin{equation*}
y_2 = \frac{m_2}{k_2}\ln\biggl[v_1\sqrt{\frac{k_2}{F_2 - m_2g}}\,\sinh\Bigl(\tfrac{t_2}{m_2}
  \sqrt{k_2(F_2 - m_2g)}\Bigr) + \cosh\Bigl(\tfrac{t_2}{m_2}\sqrt{k_2(F_2 - m_2g)}\Bigr)\biggr]
\end{equation*}
which is identical to equation~\eqref{ch4:eq:58} on page 536 of the original text---the same result
obtained for $v_1$ less than $v_{\mathrm{term}}$. This is expected, since the altitude equation does not
have a singularity at $v_1 = v_{\mathrm{term}}$ as the velocity equation does.

Add the generalization of extended Fehskens-Malewicki solution to rockets of three or more stages here.

\phantomsection
\section*{Upper Stages With Very Low Thrust}\label{supp:ch4-summary:low-thrust}

Add a section 2.2.4 which discusses upper stages with thrust equal to or less than weight. These cases are of
primarily academic interest, but they could apply to upper stages using low-thrust engines to minimize drag
and to analyses of the effect of tracking smoke (generally assumed to have negligible thrust) on the
coasting phase of flight.

Note that the functional forms of the extended Fehskens-Malewicki velocity \emph{and} altitude increment
equations change when the thrust becomes equal to the weight $(F_2 - m_2g = 0)$. Derive and present the
Fehskens-Malewicki solution for the case where thrust = weight:
\begin{equation*}
v_2 = \frac{v_1}{1 + \frac{k_2}{m_2}v_1t_2}
\end{equation*}
and
\begin{equation*}
y_2 = \frac{m_2}{k_2}\ln\biggl(1 + \frac{k_2}{m_2}v_1t_2\biggr)
\end{equation*}
Note that the ``terminal velocity'' has now become zero and the rocket's velocity decays like that of a
submarine or airship whose engines are suddenly shut down.

Present the generalization of the extended Fehskens-Malewicki equations to rockets of three or more stages.

Derive and present the extended Caporaso-Bengen solution for the case where thrust = weight:
\begin{equation*}
v_2 = \frac{v_1}{\sqrt{1 + \tfrac{2k_2}{m_2}v_1t_2}}
\end{equation*}
and
\begin{equation*}
y_2 = \frac{m_2}{k_2}\biggl(\sqrt{1 + \tfrac{2k_2}{m_2}v_1t_2} - 1\biggr)
\end{equation*}
Note that the Caporaso-Bengen solution was originally derived to provide an accurate, approximate
closed-form solution for vertical rocket flight in cases where the rocket's thrust substantially exceeds its
weight and drag. While the extended Caporaso-Bengen solution remains accurate for most practical cases of
upper stages ignited at airspeeds exceeding their terminal velocities, the accuracy of the Caporaso-Bengen
solution decreases as the thrust becomes so small that it approaches the weight because the error due to
truncation of the drag integral becomes more significant even at moderate values of $dv/dt$. For thrust
less than 1.25 times the rocket's weight the Caporaso-Bengen equations will produce acceptably accurate
results only if $v_1$ is not greater than $v_{\mathrm{fterm}}/\sqrt{2}$, where $v_{\mathrm{fterm}}$ is the
downward terminal velocity the rocket would approach in free fall (i.e., with zero thrust) if its recovery
device did not deploy:
\begin{equation*}
v_{\mathrm{fterm}} = \sqrt{\frac{m_2g}{k_2}}
\end{equation*}
Present the generalization of the extended Caporaso-Bengen equations to rockets of three or more stages.

Note that the functional forms of the extended Fehskens-Malewicki velocity and altitude increment equations
change again when the rocket's thrust is less than its weight (the case of ``boosted coasting''). Derive and
present the extended Fehskens-Malewicki solution for the case where thrust is less than weight:
\begin{equation*}
v_2 = \sqrt{\frac{m_2g - F_2}{k_2}}\,\tan\biggl[\tan^{-1}\biggl(v_1\sqrt{\frac{k_2}{m_2g - F_2}}\biggr)
  - \frac{t_2}{m_2}\sqrt{k_2(m_2g - F_2)}\biggr]
\end{equation*}
and
\begin{equation*}
y_2 = \frac{m_2}{2k_2}\ln\left[\frac{v_1^{2} + \dfrac{m_2g - F_2}{k_2}}{v_2^{2} + \dfrac{m_2g - F_2}{k_2}}
  \right]
\end{equation*}
Derive and present the extended Fehskens-Malewicki solution for the maximum altitude increment and boosted
coasting time:
\begin{equation*}
t_{\mathrm{BC}} = \frac{m_2}{\sqrt{k_2(m_2g - F_2)}}\tan^{-1}\biggl(v_1\sqrt{\frac{k_2}{m_2g - F_2}}\biggr)
\end{equation*}
and
\begin{equation*}
y_{\mathrm{BC}} = \frac{m_2}{2k_2}\ln\biggl(\frac{k_2v_1^{2}}{m_2g - F_2} + 1\biggr)
\end{equation*}
Present the generalization of the extended Fehskens-Malewicki equations to rockets of three or more stages.

Derive and present the extended Caporaso-Bengen solution for the case where thrust is less than weight:
\begin{equation*}
v_2 = \frac{m_2v_1 - (m_2g - F_2)t_2}{\sqrt{m_2^{2} + 2k_2t_2\bigl[m_2v_1 - \tfrac{t_2}{2}(m_2g - F_2)\bigr]}}
\end{equation*}
and
\begin{equation*}
y_2 = \frac{-m_2 + \sqrt{m_2^{2} + 2k_2t_2\bigl[m_2v_1 - \tfrac{t_2}{2}(m_2g - F_2)\bigr]}}{k_2}
\end{equation*}
Derive and present the extended Caporaso-Bengen solution for the maximum altitude increment and the boosted
coasting time:
\begin{gather*}
t_{\mathrm{BC}} = \frac{m_2v_1}{m_2g - F_2}
\end{gather*}
and
\begin{equation*}
y_{\mathrm{BC}} = \frac{m_2}{k_2}\biggl(\sqrt{\frac{k_2v_1^{2}}{m_2g - F_2} + 1} - 1\biggr)
\end{equation*}
Note that the Caporaso-Bengen solution for this case predicts the same ``boosted'' coasting time as would
occur with zero drag, and an altitude increment greater than that which will actually be attained but less
than that which would occur with zero drag. Note again that acceptable accuracy can be obtained from the
Caporaso-Bengen solution for this case only if $v_1$ is not greater than $v_{\mathrm{fterm}}/\sqrt{2}$.

\phantomsection
\section*{Solutions for a Non-Oscillating Rocket in the Coasting Phase}\label{supp:ch4-summary:coasting}

Add a treatment of the extended Caporaso-Bengen solution for the maximum altitude increment and coasting
time:
\begin{gather*}
t_c = \frac{v_b}{g}
\end{gather*}
and
\begin{equation*}
y_c = \frac{m_b}{k}\biggl(\sqrt{\frac{kv_b^{2}}{m_bg} + 1} - 1\biggr)
\end{equation*}
Note that acceptable accuracy can be obtained from the Caporaso-Bengen equations for this case only if $v_b$
is not greater than $v_{\mathrm{fterm}}/\sqrt{2}$.
